\documentclass[10pt]{article}
\providecommand{\FIGDIR}{fig}
\newcommand{\DIGESTDATE}{29 August 2026}
% ===================== PM Research Watch - shared preamble =====================
\usepackage[T1]{fontenc}
\usepackage[utf8]{inputenc}
\usepackage{amsmath,amssymb}
\usepackage{tgpagella}
\usepackage{tgheros}
\renewcommand{\sfdefault}{qhv}
\renewcommand{\ttdefault}{lmtt}
\usepackage{microtype}
\usepackage[english]{babel}

\usepackage[a4paper,top=17mm,bottom=17mm,left=17mm,right=17mm,
            headsep=6mm,footskip=10mm]{geometry}
\usepackage{graphicx}
\usepackage{xcolor}
\usepackage{colortbl}
\usepackage{tikz}
\usetikzlibrary{calc}
\usepackage[most]{tcolorbox}
\usepackage{titlesec}
\usepackage{fancyhdr}
\usepackage{booktabs}
\usepackage{array}
\usepackage{longtable}
\usepackage{enumitem}
\usepackage{multicol}
\usepackage{ragged2e}
\usepackage[hidelinks]{hyperref}

% ------------------------------------------------------------------ palette
\definecolor{Ink}    {HTML}{14262E}
\definecolor{Deep}   {HTML}{0F4C5C}
\definecolor{Teal}   {HTML}{1F7A8C}
\definecolor{Sky}    {HTML}{4FA3C4}
\definecolor{Sage}   {HTML}{6FA287}
\definecolor{Amber}  {HTML}{E3A008}
\definecolor{Clay}   {HTML}{B5643C}
\definecolor{Coral}  {HTML}{D1495B}
\definecolor{Violet} {HTML}{6A5B9E}
\definecolor{Slate}  {HTML}{4A6572}
\definecolor{Paper}  {HTML}{FBF9F5}
\definecolor{Mist}   {HTML}{EEEAE2}
\definecolor{Rule}   {HTML}{D9D3C9}

\pagecolor{Paper}
\color{Ink}

\hypersetup{colorlinks=true, linkcolor=Deep, urlcolor=Teal, citecolor=Teal,
            pdftitle={PM Research Watch}, pdfauthor={Automated daily digest}}

% ------------------------------------------------------------------ headings
\titleformat{\section}
  {\sffamily\bfseries\large\color{Deep}}{}{0pt}
  {}[{\vspace{1.2mm}\color{Amber}\titlerule[1.1pt]}]
\titlespacing*{\section}{0pt}{5.5mm}{2.8mm}

\titleformat{\subsection}
  {\sffamily\bfseries\normalsize\color{Teal}}{}{0pt}{}
\titlespacing*{\subsection}{0pt}{3.6mm}{1.4mm}

\setlist[itemize]{leftmargin=4.2mm, itemsep=1.1mm, topsep=1.3mm, parsep=0pt,
                  label=\textcolor{Teal}{\footnotesize$\blacktriangleright$}}

% ------------------------------------------------------------------ header/footer
\pagestyle{fancy}
\fancyhf{}
\renewcommand{\headrulewidth}{0pt}
\renewcommand{\footrulewidth}{0pt}
\fancyhead[L]{\sffamily\scriptsize\color{Slate}PM RESEARCH WATCH\;\textcolor{Rule}{\textbar}\;\DIGESTDATE}
\fancyhead[R]{\sffamily\scriptsize\color{Slate}Particulate matter monitoring \& health effects}
\fancyfoot[C]{\sffamily\scriptsize\color{Slate}\thepage}

% ------------------------------------------------------------------ components
% metric tile
\newcommand{\metric}[3]{%
  \begin{tcolorbox}[enhanced, colback=white, colframe=Rule, boxrule=0.5pt,
    arc=1.6mm, left=2.2mm, right=2.2mm, top=1.8mm, bottom=1.8mm,
    borderline west={1.7mm}{0pt}{#3}, nobeforeafter, valign=center,
    height=17mm, width=\linewidth]
    \RaggedRight
    {\sffamily\bfseries\fontsize{17}{17}\selectfont\color{#3}#1}\\[0.5mm]
    {\sffamily\fontsize{6.5}{7.4}\selectfont\color{Slate}#2}
  \end{tcolorbox}}

% pull-quote / key insight box
\newtcolorbox{keybox}[1][]{enhanced, colback=Mist, colframe=Mist,
  arc=1.6mm, left=3.2mm, right=3.2mm, top=2.4mm, bottom=2.4mm,
  borderline west={1.7mm}{0pt}{Amber}, #1}

% subtopic band
\newcommand{\band}[2]{%
  \vspace{3.2mm}\noindent
  \begin{tcolorbox}[enhanced, colback=#2, colframe=#2, arc=1.2mm,
    left=2.6mm, right=2.6mm, top=1.1mm, bottom=1.1mm, boxrule=0pt,
    width=\linewidth, nobeforeafter]
    {\sffamily\bfseries\small\color{white}#1}
  \end{tcolorbox}\vspace{1.6mm}}

% one paper entry: {short cite}{title}{journal}{doi}{body}
\newcommand{\paperentry}[5]{%
  \par\vspace{1.9mm}\noindent
  {\sffamily\bfseries\footnotesize\color{Ink}#1}\;%
  {\sffamily\fontsize{7}{8}\selectfont\color{Slate}#3}\\[0.4mm]
  {\footnotesize\itshape\color{Deep}#2}\\[0.6mm]
  {\footnotesize #4}\\[0.3mm]
  {\fontsize{6.6}{7.6}\selectfont\color{Slate}\texttt{doi:}\,\href{https://doi.org/#5}{\texttt{\detokenize{#5}}}}
  \par}

% caption style
\newcommand{\figcap}[1]{\par\vspace{1.0mm}
  {\fontsize{7.2}{8.6}\selectfont\color{Slate}\textsf{\textbf{Fig.}}\; #1\par}\vspace{1.5mm}}

\providecommand{\permil}{\textperthousand}

\setlength{\parindent}{0pt}
\setlength{\parskip}{0pt}
\renewcommand{\arraystretch}{1.18}


\begin{document}

% ==================================================================== MASTHEAD
\thispagestyle{empty}
\begin{tikzpicture}[remember picture, overlay]
  \fill[Deep] (current page.north west) rectangle
      ($(current page.north east)+(0,-67mm)$);
  \fill[Amber] ($(current page.north west)+(0,-67mm)$) rectangle
      ($(current page.north east)+(0,-68.4mm)$);
  \foreach \x/\y/\r in {22/12/0.9, 41/26/1.5, 63/9/0.7, 88/31/1.1, 112/17/1.9,
                        138/29/0.8, 159/13/1.3, 178/54/1.0, 196/21/1.6,
                        30/44/0.6, 74/57/1.2, 125/48/0.9, 168/6/0.8, 191/59/0.7,
                        52/18/1.0, 100/60/0.6, 148/20/1.1, 66/49/1.4, 205/45/0.8}
    \fill[Sky, opacity=0.30] ($(current page.north west)+(\x mm,-\y mm)$) circle (\r mm);
\end{tikzpicture}

\vspace*{4mm}
{\sffamily\fontsize{7.6}{9}\selectfont\color{Sky}%
 AUTOMATED DAILY LITERATURE DIGEST\;\textbullet\;PhD RESEARCH WATCH\par}
\vspace{2.2mm}
{\sffamily\bfseries\fontsize{27}{29}\selectfont\color{white}%
 Particulate Matter\par}
{\sffamily\fontsize{27}{29}\selectfont\color{Sky}%
 Monitoring \& Health Effects\par}
\vspace{5.0mm}
{\sffamily\fontsize{8.4}{10}\selectfont\color{white}%
 Issue 2026-08-29\;\;\textcolor{Amber}{\textbar}\;\;Entry window 29 August 2026\;\;%
 \textcolor{Amber}{\textbar}\;\;8 records screened in scope\par}
\vspace{18mm}

% -------------------------------------------------------------- metric strip
\noindent
\begin{minipage}[t]{0.190\linewidth}\metric{8}{RECORDS IN SCOPE\\(22 EXCLUDED)}{Deep}\end{minipage}\hfill
\begin{minipage}[t]{0.190\linewidth}\metric{7}{SUBTOPIC CLUSTERS\\OCCUPIED (OF 10)}{Teal}\end{minipage}\hfill
\begin{minipage}[t]{0.190\linewidth}\metric{4}{REVIEWS /\\POSITION STATEMENTS}{Sky}\end{minipage}\hfill
\begin{minipage}[t]{0.190\linewidth}\metric{5}{EFFECT ESTIMATES\\FROM 1 STUDY}{Amber}\end{minipage}\hfill
\begin{minipage}[t]{0.190\linewidth}\metric{2}{RECORDS FROM\\SUB-SAHARAN AFRICA}{Clay}\end{minipage}

\vspace{6mm}

% -------------------------------------------------------------- signal of day
\section*{\sffamily\bfseries\color{Deep}Signal of the day}
\vspace{-1.5mm}{\color{Amber}\rule{\linewidth}{1.1pt}}\vspace{2.5mm}

\begin{keybox}
{\sffamily\bfseries\small\color{Deep}Thyroid function joins the list of endpoints with a
coherent air-pollution dose--response --- and PM$_{10}$, not PM$_{2.5}$, carries the
strongest association.}\\[1.4mm]
\footnotesize
Park et al.\ (\textit{J Endocrinol Invest}) screened \textbf{16{,}193 Korean adults} with
TSH in the assayable 0.1--10.0\,mIU/L range against CMAQ-modelled annual PM$_{2.5}$,
PM$_{10}$, NO$_2$, SO$_2$ and CO. Every pollutant is positively associated with TSH and
negatively with free thyroxine ($P<0.001$), and for the combined
\textbf{high-TSH/low-FT4 phenotype} the adjusted odds ratios per IQR order as
\textbf{PM$_{10}$ 1.55 (95\% CI 1.42--1.69)} $>$ SO$_2$ 1.38 (1.28--1.49) $>$ NO$_2$ 1.36
(1.23--1.50) $>$ \textbf{PM$_{2.5}$ 1.25 (1.17--1.34)} $>$ CO 1.23 (1.15--1.31), with
stronger effects in women. That PM$_{10}$ outranks PM$_{2.5}$ is worth sitting with: five
collinear pollutants from a single CMAQ field, entered in separate single-pollutant
models, cannot separate coarse-fraction mechanism from correlation structure --- but if
the coarse fraction really is doing the work, it is the fraction most low-cost optical
networks measure worst, because PM$_{10}$ sizing in a nephelometric unit is the least
reliable channel it has. Cross-sectional design and a single-hospital screening
population cap the inference at hypothesis strength.
\end{keybox}

\vspace{3mm}

% -------------------------------------------------------------- exec brief
\section*{\sffamily\bfseries\color{Deep}What changed today}
\vspace{-1.5mm}{\color{Amber}\rule{\linewidth}{1.1pt}}\vspace{2.5mm}

\footnotesize
\begin{itemize}
\item \textbf{A thin day, and half of it is synthesis.} Eight records over seven clusters,
of which \textbf{four are reviews or position statements}. No cluster holds more than two.

\item \textbf{Sub-Saharan Africa appears twice after an eight-issue absence.} A West Rand
time-activity survey (330 households near gold-mine tailings) and a meningitis-belt dust
review. The time-activity paper is the more useful for exposure work: it reports indoor
fractions \textbf{roughly 10 percentage points below the USEPA Exposure Factors Handbook
values} that most exposure models silently import, which biases indoor--outdoor
apportionment for any African deployment using default time budgets.

\item \textbf{Microplastics enter through a particle-toxicology route.} Chen et al.\
(\textit{Part Fibre Toxicol}) quantify placental microplastic burden by pyrolysis-GC-MS
and stratify transcriptome, proteome and metabolome by it --- the first record in this
watch to treat placental MP burden as a graded internal dose rather than a
presence/absence finding.

\item \textbf{A cautionary machine-learning result.} Yu et al.\ (\textit{Indoor Air})
predict seven self-reported symptoms in Milan dormitories from continuous environmental
quality data: \textbf{mean AUC 0.595 (0.398--0.705)}, only four models above 0.6, from 74
occupant samples and 131 candidate predictors. The paper's own observation is the
interesting one --- dominant predictors were \textbf{outdoor statistics and
distributional shape features, not mean indoor concentrations}. Mean-value summaries of a
sensor time series discard most of the health-relevant information.

\item \textbf{Two records that are mostly agenda-setting.} The SynAir-G statement on
children's indoor air quality argues for chemical$\times$biological synergy as the
policy-relevant unit of exposure; the poultry-worker airway review restates a
well-established organic-dust hazard without new quantification.

\item \textbf{Rejected but worth naming.} Two directly on-topic Elsevier deposits --- a
hybrid AQI forecasting framework (\textit{Atmos Pollut Res}) and an operational
ragweed-pollen forecast evaluation (\textit{Environ Pollut}) --- arrived with no abstract
in any indexed source and are carried to the retry list rather than summarised from
titles.
\end{itemize}

\clearpage

% ==================================================================== ANALYTICS
\section*{\sffamily\bfseries\color{Deep}Corpus at a glance}
\vspace{-1.5mm}{\color{Amber}\rule{\linewidth}{1.1pt}}\vspace{3mm}

\noindent\includegraphics[width=\linewidth]{\FIGDIR/f1_subtopics.png}
\figcap{Records per subtopic cluster: 8 records over \textbf{seven of ten} clusters, with
\textit{Occupational \& indoor} the only cluster holding 2. \textit{Cardiovascular \&
metabolic}, \textit{Neuro / mental health} and \textit{Reproductive \& developmental} are
empty --- although the placental microplastics record is reproductive in substance and
sits under \textit{Mechanistic toxicology} because its contribution is molecular rather
than obstetric. A flat distribution like this one carries no signal about the field; it
is what a low-volume deposit day looks like.}

\vspace{2mm}
\noindent\includegraphics[width=\linewidth]{\FIGDIR/f2_design_endpoint.png}
\figcap{Left --- study architecture, and the dominant fact of this issue:
\textbf{\textit{Review / synthesis} is 4 of 8}, against \textit{Observational ---
cross-sectional} 2, \textit{Experimental / toxicology} 1 and \textit{Measurement
campaign} 1. A day with more synthesis than primary data is a deposit-calendar artefact,
not a statement about research direction. Right --- health endpoint: one record (the
West Rand time-activity survey) carries none, and no endpoint repeats.}

\vspace{2mm}
\noindent\includegraphics[width=\linewidth]{\FIGDIR/f3_metric_geography.png}
\figcap{Left --- particle metric. \textit{PM (unspeciated) / emissions} 4 dominates,
because half the corpus characterises a mixed or organic dust rather than a regulated
mass fraction; one each for \textit{PM$_{2.5}$ only},
\textit{PM$_{2.5}$ + PM$_{10}$ jointly}, \textit{composition / speciation} (placental
polymer burden) and \textit{bioaerosol / coarse}. Right --- geography: Europe 2,
\textbf{Sub-Saharan Africa 2}, multi-region 2, China 1, East Asia ex-China 1. The African
pair ends an eight-issue absence running from 21 August; the last previous African
record was the Addis Ababa study of 20 August.}

\vspace{2mm}
\noindent\includegraphics[width=\linewidth]{\FIGDIR/f4_heatmap.png}
\figcap{Subtopic $\times$ architecture cross-tabulation. With 8 records over seven
clusters and four architectures, \textbf{only one cell holds 2} ---
\textit{Occupational \& indoor} $\times$ \textit{Review / synthesis} --- and the other six
occupied cells hold 1. There is no structure to read here; the panel is carried for
continuity of the series.}

\vspace{2mm}
\noindent\includegraphics[width=\linewidth]{\FIGDIR/f5_forest.png}
\figcap{Five estimates, \textbf{all from a single study} (Park et al.), all for the same
outcome and all on a common per-IQR contrast, so this is one of the few panels in this
watch where a pollutant-versus-pollutant reading is internally consistent. It is
\emph{not} a potency comparison in any causal sense: the five exposures come from one
CMAQ field, are mutually collinear, and were entered in separate single-pollutant models,
so the ordering (PM$_{10}$ $>$ SO$_2$ $>$ NO$_2$ $>$ PM$_{2.5}$ $>$ CO) partly reflects
each pollutant's IQR width and its correlation with the others rather than an independent
effect. The intervals are narrow because $n$ is large, not because confounding is small.}

\vspace{2mm}
\noindent\includegraphics[width=\linewidth]{\FIGDIR/f6_lifecourse.png}
\figcap{Life-course windows. \textbf{Four records across three of five windows}: the
in-utero cell holds the placental microplastics study, the childhood cell the SynAir-G
schools statement, and the working-age cell both the Korean screening cohort and the
Milan dormitory panel. Puberty/adolescence and 65+ are empty.}

\clearpage

% ==================================================================== DIGEST
\section*{\sffamily\bfseries\color{Deep}Paper-level digest}
\vspace{-1.5mm}{\color{Amber}\rule{\linewidth}{1.1pt}}\vspace{1mm}

{\fontsize{7.4}{9}\selectfont\color{Slate}Ordered by cluster; the 29 August entry window
only.\par}

\band{Sensing, forecasting \& instrumentation\hfill 1 record}{Amber}

\paperentry{Yu et al.\ 2026b}
{Machine learning-based prediction of occupant symptom risk using indoor and outdoor
environmental quality data in university dormitories}
{Indoor Air}
{Seven monitoring periods across two Milan student residences, May 2024 to June 2025,
pairing continuous indoor and outdoor environmental quality data with occupant symptom
surveys. From \textbf{74 occupant samples} and 131 candidate predictors (reduced per
symptom to 3--40 by permutation importance), seven binary symptom classifiers were fitted
under leave-one-period-out cross-validation. Performance is weak and honestly reported:
\textbf{mean AUC 0.595, range 0.398--0.705}, four of seven above 0.6; the higher combined
accuracy figures appear only after dropping the two least specific symptoms. With $n=74$
against 131 candidate predictors the feature-selection step is almost certainly
overfitting even under LOPO-CV. The transferable finding is the feature ranking:
\textbf{outdoor statistics and distributional shape features beat mean indoor
concentrations}, which argues directly against summarising a sensor time series by its
mean.}
{10.1155/ina/8230467}

\band{Exposure assessment \& modelling\hfill 1 record}{Amber}

\paperentry{Kalumbi et al.\ 2026}
{Age-specific time-spent patterns among residents living near gold mine tailings storage
facilities in the West Rand, South Africa}
{J Expo Sci Environ Epidemiol}
{Cross-sectional structured-questionnaire survey of \textbf{330 households in Kagiso},
West Rand, resolving indoor and outdoor time across five age categories with
Kruskal--Wallis tests and Bonferroni-adjusted Dunn's comparisons. Indoor time runs
\textbf{66.6--84.8\% of a weekday and 76.0--82.7\% of a weekend day}
(18.28$\pm$3.02\,h weekdays, 19.00$\pm$1.99\,h weekends); young children and older adults
are indoors most, and weekend indoor time rises significantly for children, adolescents
and adults but not for the two extreme age groups. The headline for exposure modelling is
that indoor time is \textbf{nearly 10 percentage points below USEPA Exposure Factors
Handbook defaults}. Self-reported recall with no diary or wearable validation, and a
single district, so the magnitude is soft --- the direction of the bias in imported
default time budgets is the usable result.}
{10.1038/s41370-026-00972-6}

\band{Mechanistic toxicology\hfill 1 record}{Violet}

\paperentry{Chen et al.\ 2026j}
{Multi-omics integration reveals human placental molecular signatures associated with
microplastics exposure}
{Part Fibre Toxicol}
{Placental microplastic burden quantified by \textbf{pyrolysis-gas chromatography--mass
spectrometry}, samples stratified into high- and low-burden groups, then transcriptomic,
proteomic and metabolomic profiling integrated through the DIABLO framework with network
analysis for regulatory hubs. High-burden placentae show multi-layer molecular
perturbation centred on immune dysregulation. The design's strength is that MP burden is
treated as a \emph{graded internal dose} rather than a detection/non-detection flag,
which separates it from most human-tissue microplastics work. Its limits are the usual
ones for cross-sectional human tissue: burden is measured once at delivery, exposure
timing is unknown, the omics contrasts are correlational, no sample size appears in the
abstract, and Py-GC-MS quantifies polymer \emph{mass} without particle size or count ---
so the dose metric is mass, not number or surface area.}
{10.1186/s12989-026-00697-2}

\clearpage

\band{Occupational \& indoor\hfill 2 records}{Clay}

\paperentry{Ramirez et al.\ 2026}
{Airway diseases and occupational exposures in poultry industry workers: a comprehensive
review}
{Curr Allergy Asthma Rep}
{Review of occupational asthma and chronic obstructive lung disease in poultry workers,
attributing risk to the barn's mixed aerosol --- feather and feed dust, mites, moulds,
ammonia and bacterial endotoxin --- acting through airway inflammation and
hypersensitivity, with a symptom spectrum from wheeze and chest tightness to chronic
cough and dyspnoea. Emphasis falls on early identification through occupational
history-taking and on barn control practices. \textbf{No prevalence figures, no
exposure--response estimates and no search or inclusion strategy} are reported, so this is
a clinical orientation piece rather than an evidence synthesis. Its relevance here is the
exposure architecture: an occupational aerosol whose health-relevant fraction is defined
by biological composition, which a mass-only monitor cannot resolve at all.}
{10.1007/s11882-026-01289-y}

\paperentry{Siemens et al.\ 2026}
{A SynAir-G position statement for the improvement of indoor air quality for children}
{Indoor Air}
{Position statement from the EU-funded SynAir-G project (18 partners, 11 countries) within
the seven-project IDEAL Cluster, arguing that indoor air has been under-regulated relative
to outdoor air despite most exposure occurring indoors, and that the policy-relevant unit
is the \textbf{synergistic interaction between chemical and biological pollutants} rather
than any single-pollutant concentration. Schools and children are the target setting. As a
consensus document it presents no new measurements, no systematic evidence appraisal and
no quantified synergy, so the central claim is asserted rather than demonstrated. The
instrumentation implication is nonetheless concrete: a synergy framing requires co-located
chemical and bioaerosol sensing at classroom scale, which no current low-cost platform
delivers.}
{10.1155/ina/9690417}

\band{Respiratory \& allergic\hfill 1 record}{Sage}

\paperentry{Silva et al.\ 2026}
{Shared mechanisms linking PM$_{2.5}$ exposure, obstructive sleep apnea and obesity}
{Front Sleep}
{Narrative review of the three-way relationship between chronic PM$_{2.5}$ exposure,
obstructive sleep apnea and obesity, proposing systemic inflammation, oxidative stress,
neuroendocrine dysregulation and gut microbiota alteration as shared pathways, and noting
that intermittent hypoxia and sleep fragmentation feed back into weight gain. The
mechanistic overlap is plausible and the review is explicit that it is
hypothesis-organising, but \textbf{no effect estimates, no study inventory and no
appraisal of confounding by adiposity} appear --- and adiposity confounds both the
exposure (residential location) and the outcome. Read alongside the CHARLS sleep-duration
result carried on 28 August, the practical gap is identical in both: no objective sleep
measurement anywhere in this literature.}
{10.3389/frsle.2026.1918405}

\clearpage

\band{Other clinical endpoints\hfill 1 record}{Slate}

\paperentry{Park et al.\ 2026b}
{Association between annual air pollutant exposure and thyroid function: a retrospective
cross-sectional study}
{J Endocrinol Invest}
{\textbf{16{,}193 adults} (8{,}812 male, 7{,}381 female; mean age 53.8$\pm$13.1\,y)
screened at Seoul National University with TSH within 0.1--10.0\,mIU/L, against
CMAQ-estimated annual PM$_{2.5}$, PM$_{10}$, NO$_2$, SO$_2$ and CO. All five pollutants
associate positively with TSH and negatively with free thyroxine ($P<0.001$), with
significant non-linearity for PM$_{2.5}$ and NO$_2$ against log TSH. For the combined
high-TSH/low-FT4 phenotype the per-IQR adjusted odds ratios run
\textbf{PM$_{10}$ 1.55 (1.42--1.69)}, SO$_2$ 1.38 (1.28--1.49), NO$_2$ 1.36 (1.23--1.50),
\textbf{PM$_{2.5}$ 1.25 (1.17--1.34)} and CO 1.23 (1.15--1.31), stronger among women. The
binding limitations: exposure is a modelled annual field at residence with no personal
component, the design is cross-sectional so temporality is unresolved, and the five
pollutants are collinear by construction --- the ordering must not be read as
differential potency.}
{10.1007/s40618-026-02965-6}

\band{Burden, policy \& mitigation\hfill 1 record}{Coral}

\paperentry{Cliff et al.\ 2026}
{How do climatic factors directly influence the incidence and risk of meningococcal
meningitis across the African meningitis belt?}
{GeoHealth}
{Narrative review over the 26-country African meningitis belt, where outbreaks cluster
strongly in the dry season, examining dust, wind speed, temperature, rainfall and land
cover as candidate drivers. \textbf{Atmospheric dust and wind speed carry the strongest
statistical association and the greatest predictive probability} of the variables
examined, with mucosal damage from inhaled mineral dust the standing mechanistic
hypothesis. The review's own central caveat is the right one: these climatic variables are
severely multicollinear within the dry season, so attributing incidence to dust rather
than to temperature or humidity is not identified by the observational record. Carried
here because it is one of the few live health endpoints where the \emph{coarse} mineral
fraction, not PM$_{2.5}$, is the hypothesised agent --- and where satellite dust products
substitute for a ground network that does not exist.}
{10.1029/2025gh001749}

\clearpage

% ==================================================================== TABLE
\section*{\sffamily\bfseries\color{Deep}Corpus register}
\vspace{-1.5mm}{\color{Amber}\rule{\linewidth}{1.1pt}}\vspace{2.5mm}

{\fontsize{6.7}{8.4}\selectfont
\setlength{\tabcolsep}{2.4pt}
\begin{longtable}{@{}>{\RaggedRight}p{25mm}>{\RaggedRight}p{41mm}>{\RaggedRight}p{27mm}%
>{\RaggedRight}p{23mm}>{\RaggedRight}p{18mm}>{\RaggedRight}p{15mm}@{}}
\toprule
\textsf{\textbf{First author}} & \textsf{\textbf{Focus}} & \textsf{\textbf{Journal}} &
\textsf{\textbf{Design}} & \textsf{\textbf{Metric}} & \textsf{\textbf{Region}}\\
\midrule
\endfirsthead
\toprule
\textsf{\textbf{First author}} & \textsf{\textbf{Focus}} & \textsf{\textbf{Journal}} &
\textsf{\textbf{Design}} & \textsf{\textbf{Metric}} & \textsf{\textbf{Region}}\\
\midrule
\endhead
\multicolumn{6}{@{}l@{}}{\cellcolor{Mist}\textsf{\textbf{\textcolor{Deep}{Neuro / mental health}}}}\\[0.6mm]
\href{https://doi.org/10.1001/jamanetworkopen.2026.25015}{Dales} & Wildfire-sourced PM$_{2.5}$ and ED visits for migraine and other primary headache & \textit{JAMA Netw Open} & Case-crossover & PM$_{2.5}$ (wildfire) & Canada\\
\href{https://doi.org/10.1093/ageing/afag215}{Glover} & Umbrella meta-review of urban environmental exposures and cognitive health & \textit{Age Ageing} & Umbrella review / meta-analysis & PM$_{2.5}$ & Global\\
\href{https://doi.org/10.1016/j.nsa.2026.107019}{Goossens} & In utero air pollution exposure and pre-/postnatal brain development & \textit{Neurosci Appl} & Narrative review & PM$_{2.5}$ / PM$_{10}$ / PM$_{\mathrm{coarse}}$ & Global\\
\href{https://doi.org/10.1016/j.isci.2026.116617}{Ma} & PM$_{2.5}$-induced neurotoxicity and depression: exposure risks and mechanisms & \textit{iScience} & Narrative review & PM$_{2.5}$ + components & Global\\
\href{https://doi.org/10.3760/cma.j.cn112150-20250804-00752}{Pan} & Maternal PM exposure and offspring neurodevelopment: effects and mechanisms & \textit{Zhonghua Yu Fang Yi Xue Za Zhi} & Narrative review & PM (bulk) & China\\
\addlinespace[0.9mm]
\multicolumn{6}{@{}l@{}}{\cellcolor{Mist}\textsf{\textbf{\textcolor{Teal}{Cardiovascular \& metabolic}}}}\\[0.6mm]
\href{https://doi.org/10.3760/cma.j.cn112338-20251120-00843}{Huang} & PM$_{2.5}$ and dyslipidaemia in rural Chinese children; dietary effect modification & \textit{Zhonghua Liu Xing Bing Xue Za Zhi} & Cross-sectional (multistage) & PM$_{2.5}$ & China\\
\href{https://doi.org/10.32394/pe/218966}{Kazmierska} & Air pollution and cardiovascular disease: review of current data & \textit{Przegl Epidemiol} & Narrative review & PM$_{2.5}$ / PM$_{10}$ & Global\\
\href{https://doi.org/10.1007/s00125-026-06814-2}{Xi} & Air pollution, built environment and temperature vs. prediabetes phenotypes & \textit{Diabetologia} & Prospective cohort & PM$_{2.5}$ & Germany\\
\href{https://doi.org/10.1161/JAHA.125.046917}{Yang} & Short-term PM$_{2.5}$ and hospitalisation/ED use in atrial fibrillation & \textit{J Am Heart Assoc} & Case-crossover & PM$_{2.5}$ (wildfire-influenced) & USA\\
\addlinespace[0.9mm]
\multicolumn{6}{@{}l@{}}{\cellcolor{Mist}\textsf{\textbf{\textcolor{Sky}{Reproductive \& developmental}}}}\\[0.6mm]
\href{https://doi.org/10.1016/j.jes.2025.10.043}{Chen} & Prenatal air pollution, fetal growth trajectory phenotypes and child neurodevelopment & \textit{J Environ Sci (China)} & Prospective cohort & PM$_{2.5}$ / PM$_{10}$ & China\\
\href{https://doi.org/10.1016/j.ecoenv.2026.120566}{Huang} & Joint prenatal bisphenol + air-pollutant mixture exposure and preschool IQ & \textit{Ecotoxicol Environ Saf} & Prospective cohort (mixtures) & PM$_{2.5}$ / PM$_{10}$ & China\\
\href{https://doi.org/10.1007/s40572-026-00552-8}{Peterson} & Ambient air pollution, PAHs and pubertal development: critical review & \textit{Curr Environ Health Rep} & Critical review & PM$_{2.5}$ components & Global\\
\href{https://doi.org/10.1016/j.envres.2026.125313}{Xie} & Long-term PM$_{2.5}$ and timing of pubertal development in boys and girls & \textit{Environ Res} & Prospective cohort & PM$_{2.5}$ & China\\
\addlinespace[0.9mm]
\multicolumn{6}{@{}l@{}}{\cellcolor{Mist}\textsf{\textbf{\textcolor{Sage}{Respiratory \& allergic}}}}\\[0.6mm]
\href{https://doi.org/10.1007/s00484-026-03289-z}{Bayar Muluk} & Climate change and otolaryngology: airway disease impacts and emission reduction & \textit{Int J Biometeorol} & Narrative review & PM$_{2.5}$ & Global\\
\href{https://doi.org/10.4103/mgr.MEDGASRES-D-25-00164}{Liao} & Short-term air pollution, respiratory symptoms and lung function in asthma & \textit{Med Gas Res} & Panel study & PM$_{2.5}$ / PM$_{10}$ & China\\
\href{https://doi.org/10.3390/epidemiologia7040097}{Sanchez de Prada} & Ambient PM and asthma-coded admissions in children 0-3 y, Spain & \textit{Epidemiologia} & Ecological panel & PM$_{2.5}$ / PM$_{10}$ & Spain\\
\href{https://doi.org/10.1016/j.jhazmat.2026.143076}{Yang} & Allergenic vs. toxic risk of airborne seed hairs compared with ambient PM$_{2.5}$ & \textit{J Hazard Mater} & Field measurement & PM$_{2.5}$ + bioaerosol & China\\
\addlinespace[0.9mm]
\multicolumn{6}{@{}l@{}}{\cellcolor{Mist}\textsf{\textbf{\textcolor{Violet}{Mechanistic toxicology}}}}\\[0.6mm]
\href{https://doi.org/10.1016/j.jes.2025.12.074}{Ding} & PM$_{2.5}$ promotes vascular calcification via miR-27b-5p/HES1 signalling & \textit{J Environ Sci (China)} & In vitro + organoid + murine & PM$_{2.5}$ & China\\
\href{https://doi.org/10.1007/s00360-026-01698-5}{Maraschi} & Osmoregulatory disruption and metal bioaccumulation from settleable atmospheric PM & \textit{J Comp Physiol B} & Controlled exposure (ecotox) & Settleable APM & Brazil\\
\href{https://doi.org/10.1016/j.jes.2025.11.018}{Wu} & Macrophage NLRP3 inflammasome in real-ambient PM-induced hepatic glucose dysregulation & \textit{J Environ Sci (China)} & Real-ambient murine + co-culture & Real-ambient PM & China\\
\addlinespace[0.9mm]
\multicolumn{6}{@{}l@{}}{\cellcolor{Mist}\textsf{\textbf{\textcolor{Clay}{Occupational \& indoor}}}}\\[0.6mm]
\href{https://doi.org/10.3389/fmedt.2026.1796828}{Chen} & Nested suction instrument to cut intraoperative surgical-smoke PM$_{2.5}$ & \textit{Front Med Technol} & Bench / device experiment & PM$_{2.5}$ (surgical smoke) & China\\
\href{https://doi.org/10.1093/annweh/wxag055}{Cole} & Determinants of respirable crystalline silica exposure in Australian tunnelling & \textit{Ann Work Expo Health} & Observational exposure study & Respirable dust / RCS & Australia\\
\href{https://doi.org/10.1016/j.jhazmat.2026.143081}{Ding} & Stage-resolved hair Zn-isotope dynamics as a firefighting combustion-exposure tracer & \textit{J Hazard Mater} & Repeated-measures biomarker & Combustion-derived PM & China\\
\href{https://doi.org/10.3390/jox16040119}{Ramsden} & Xenobiotic hazards in aircraft cabin air & \textit{J Xenobiot} & Review + risk economics & Ultrafine particles & Global\\
\addlinespace[0.9mm]
\multicolumn{6}{@{}l@{}}{\cellcolor{Mist}\textsf{\textbf{\textcolor{Amber}{Sensing, forecasting \& instrumentation}}}}\\[0.6mm]
\href{https://doi.org/10.1016/j.jenvman.2026.130515}{Shi} & Adaptive spatiotemporal graph transformer for multi-site multi-step PM$_{2.5}$ forecasting & \textit{J Environ Manage} & Deep-learning model & PM$_{2.5}$ & China / India\\
\href{https://doi.org/10.3390/bios16070380}{Tsuchiya} & Seasonal and indoor-outdoor characteristics of fluorescent aerosol particles in offices & \textit{Biosensors} & Real-time sensor campaign & Bioaerosol / coarse & Japan\\
\addlinespace[0.9mm]
\multicolumn{6}{@{}l@{}}{\cellcolor{Mist}\textsf{\textbf{\textcolor{Amber}{Exposure assessment \& modelling}}}}\\[0.6mm]
\href{https://doi.org/10.1029/2025GH001621}{Solorzano-Araque} & Multiscale exposure modelling: LOTOS-EUROS + mobility surveys + particle toxicity & \textit{GeoHealth} & CTM + mobility integration & PM$_{2.5}$ / PM$_{10}$ & Colombia\\
\addlinespace[0.9mm]
\multicolumn{6}{@{}l@{}}{\cellcolor{Mist}\textsf{\textbf{\textcolor{Coral}{Burden, policy \& mitigation}}}}\\[0.6mm]
\href{https://doi.org/10.1002/hsr2.72886}{Karamzad} & Ischaemic heart disease burden in MENA 1990-2021; PM attributable fraction 28.2\% & \textit{Health Sci Rep} & GBD secondary analysis & Ambient PM & MENA\\
\href{https://doi.org/10.1371/journal.pone.0347989}{Rees} & HEPA filtration to reduce respiratory infection in care homes & \textit{PLoS One} & RCT process evaluation & Filtered PM & UK\\
\href{https://doi.org/10.1016/j.jes.2026.01.009}{Shan} & Policy-driven emission reductions and post-pandemic rebound, high-resolution inventories & \textit{J Environ Sci (China)} & Emission inventory & PM emissions & China\\
\addlinespace[0.9mm]
\multicolumn{6}{@{}l@{}}{\cellcolor{Mist}\textsf{\textbf{\textcolor{Slate}{Other clinical endpoints}}}}\\[0.6mm]
\href{https://doi.org/10.3390/dj14070440}{Lupita} & PM$_{2.5}$, preventive-care capacity and caries/edentulism burden across the EU & \textit{Dent J} & Ecological panel & PM$_{2.5}$ & EU-27\\
\href{https://doi.org/10.1016/j.jdin.2026.06.008}{Park} & Long-term PM exposure and risk of androgenetic alopecia & \textit{JAAD Int} & Nationwide cohort & PM$_{2.5}$ / PM$_{10}$ & South Korea\\
\href{https://doi.org/10.3171/2026.3.PEDS2615}{Sun} & Environmental exposure and area deprivation vs. paediatric hydrocephalus outcomes & \textit{J Neurosurg Pediatr} & Registry cohort + geospatial & PM$_{2.5}$ (satellite-derived) & USA\\
\addlinespace[0.9mm]
\bottomrule
\end{longtable}}


\vspace{4mm}
\begin{keybox}
{\sffamily\bfseries\footnotesize\color{Deep}Provenance and method}\\[1.2mm]
{\fontsize{7.2}{8.8}\selectfont
Entry window \textbf{2026-08-29 to 2026-08-29}, a single day of deposits.
\textbf{Sources.} PubMed E-utilities on \texttt{[Date - Entry]}, run as two separate
queries (health axis 20, instrumentation axis 3). Europe PMC \texttt{CREATION\_DATE}
returned 0. Crossref by-ISSN returned 79 deposits across the harvester's two-day window;
because that window opened at \texttt{last\_entry\_date + 1} while the 28 August issue was
still being written, \textbf{57 of the 79 carry a Crossref creation date of 28 August and
were screened in the previous issue} --- they are excluded here rather than re-screened,
so this issue's Crossref pool is the 22 deposits created on 29 August. OpenAlex remains
broken for entry-date windowing; arXiv contributed nothing in window.
\textbf{Screening.} \textbf{30 unique} candidates after PMID/DOI collapse and removal of
records already in \texttt{seen.json} (15 of the 23 PubMed hits were already published in
the 28 August issue), \textbf{8 retained and 22 rejected} with reason
logged to \texttt{state/rejected.jsonl}. Rejections were dominated by non-particle
environmental science (\textit{Sci Total Environ} soil, water and carbon-flux deposits
plus four retraction notices), non-airborne chemical exposure (bisphenol A, 2,4-DTBP,
endocrine-disruptor mixtures) and \textit{Meas Sci Technol} non-atmospheric metrology.
\textbf{Two in-scope deposits} --- an AQI forecasting framework and a ragweed-pollen
operational forecast --- had no retrievable abstract and were carried to the retry list
rather than summarised.
\textbf{Effect estimates.} Five estimates, all from Park et al.; the figure caption states
why the between-pollutant ordering is not a potency comparison.
\textbf{Pipeline changes made this run.} Six design labels and \textit{Italy} were added
to \texttt{plots.py} in the same edit that wrote them onto records.
\textbf{Labels.} Subtopic, design, particle-metric and geography are single-label by
dominant contribution --- a judgement, not a controlled vocabulary. All 8 DOIs passed
\texttt{check\_dois.py} (0 fail, 0 warn) before build.\par}
\end{keybox}

\end{document}
